\section{Introduction}
\label{sec:intro}

The fine-tuning of large, pre-trained foundation models (such as CLIP \cite{radford2021clip} or BERT \cite{devlin2018bert}) has become the standard paradigm for deploying highly capable models on specialized downstream tasks. However, maintaining distinct checkpoints for every individual task introduces substantial storage and deployment costs, and prevents multi-task generalist capabilities. Consequently, model merging \cite{wortsman2022model, choshen2022fusing} has emerged as an attractive, zero-inference-overhead alternative. Rather than training a unified multi-task model from scratch, model merging techniques combine independently fine-tuned expert weights directly in the parameter space to obtain a single multi-task model.

Among these, Task Arithmetic \cite{ilharco2023editing} is the foundational framework. It extracts a ``task vector'' for each expert by subtracting the pre-trained base weights from the fine-tuned expert weights, and then adds these task vectors back to the base model using a global scaling hyperparameter. While task arithmetic is incredibly simple, its performance is highly sensitive to the relative scales of the task updates. In practice, tasks with large distribution shifts or harder objectives undergo significantly larger parameter shifts during fine-tuning. This discrepancy causes their task vectors to possess disproportionately large Frobenius norms compared to other tasks. When directly summed, these high-magnitude task vectors dominate the representation space of the merged model, causing destructive interference and severely degrading performance on the other tasks.

To address this representation imbalance, state-of-the-art Test-Time Adaptation (TTA) merging techniques, such as AdaMerging \cite{yang2024adamerging} and SyMerge \cite{yu2023symerge}, have been proposed. These methods introduce task-wise or layer-wise merging coefficients which they optimize at test-time using a small calibration set of unlabeled images. The standard objective for this optimization is the minimization of joint prediction entropy across tasks. While these optimization methods have been widely adopted, we argue that this paradigm exhibits significant complexity and a lack of robustness under task-difficulty imbalances. We expose a fundamental vulnerability in this approach, which we term the \textbf{Overfitting-Optimizer Paradox}: when optimizing merging coefficients using joint entropy minimization on small calibration sets, the optimizer is naturally biased toward easy, low-entropy tasks. To minimize joint entropy, the optimization process frequently suppresses the coefficients of harder, high-entropy tasks (e.g., FashionMNIST) from the default $0.30$ down to approximately $0.23 - 0.24$, while inflating the coefficients of easy tasks (e.g., MNIST). Although these optimized coefficients remain non-zero, this relative suppression combined with easy-task gradient dominance is sufficient to degrade their representation, resulting in a delicate, transductive overfit state that generalizes poorly and actively suppresses difficult tasks. (Crucially, we note that this paradox is fundamentally a consequence of the unsupervised nature of the test-time adaptation objective, i.e., joint prediction entropy minimization on unlabeled data; under a supervised objective, this task suppression would not occur, but supervised adaptation is severely limited in practice as it demands high-quality labeled downstream data). Moreover, this paradigm introduces multiple hyperparameter choices (learning rates, batch sizes, optimizer types, and training epochs) and demands the availability of unlabeled calibration data at test time, severely limiting its applicability.

In this work, we return to the core principles of Occam's razor. We argue that the physical and geometric imbalances of task updates should be resolved analytically, rather than through computationally demanding test-time optimization loops. To this end, we introduce \textbf{Norm-Equalized Task Arithmetic (NETA)}, a training-free, parameter-free, and data-free closed-form model merging method. 

NETA operates under a simple geometric assumption: at each layer of a deep neural network, every task expert should contribute a parameter update of exactly identical magnitude. NETA dynamically scales the task vectors at each individual layer to ensure that their Frobenius norms are perfectly equalized to the layer's mean norm before merging. This transformation guarantees \textit{perfect isotropic magnitude balance} across tasks at every level of the network representation. It does so without introducing a single new hyperparameter, requiring zero calibration data, zero gradient descent steps, and zero backpropagation.

We evaluate NETA against standard Task Arithmetic, TIES-Merging \cite{yadav2023resolving}, and SOTA test-time AdaMerging baselines using a CLIP ViT-B/32 backbone across four diverse visual classification tasks (MNIST, FashionMNIST, CIFAR-10, and SVHN) across three independent random seeds. Our key contributions are:
\begin{itemize}
    \item We present \textbf{NETA}, a closed-form weight-space transformation that analytically resolves task vector magnitude imbalances, ensuring equal representation strength across tasks.
    \item We empirically demonstrate that NETA outperforms vanilla Task Arithmetic on MNIST ($+0.26\%$) and FashionMNIST ($+0.65\%$) zero-shot. While NETA's role as an isotropic regularizer slightly curtails peak performance on dominant tasks like SVHN (resulting in a marginal drop in average multi-task accuracy from $87.76\%$ to $87.17\%$), it successfully prevents representation hijacking and ensures multi-task fairness.
    \item We expose the \textbf{Overfitting-Optimizer Paradox} of test-time adaptation, showing that Task-Wise AdaMerging suffers a catastrophic drop of \textbf{-4.56\%} on FashionMNIST because the entropy optimizer favors easier tasks and suppresses harder ones. We also present a continuous \textbf{$\alpha$-relaxed NETA} framework that allows practitioners to smoothly balance peak task performance and representation fairness.
    \item We show that NETA is highly robust and stable, achieving strong generalization and competitive performance across random seeds, demonstrating that simple physical constraints can offer a robust and stable alternative to test-time optimization.
\end{itemize}
